% !PRINTED 39
\weberpage{39}
%<<CONT

\gesperrt{rationalen Functionen $F(x),\varphi(x)$ und zwar $\varphi(x)$ von niedrigerem Grade als $f(x)$ so bestimmen, dass die Gleichung}
\begin{equation}
 F(x)f(x)+\varphi(x)\psi(x)=\chi(x)
\tag{26}
\end{equation}
\gesperrt{identisch befriedigt ist.}

\begin{center}
{\large \S.\ 7.}

\vspace{0.5em}
{\large\gesperrt{Producte linearer Factoren.}}
\end{center}

Nach \S~1 erhalten wir durch Multiplication ganzer Functionen ebensolche Functionen von höherem Grade, und zwar bestimmt sich der Grad des Productes durch die Summe der Grade der einzelnen Factoren.

Wenn wir also $n$ lineare Factoren mit einander multipliciren, so entsteht eine Function $n$ten Grades, deren Bildungsweise wir etwas genauer untersuchen müssen.

Wir wollen die linearen Factoren in der einfachen Form $x-\alpha_1$, $x-\alpha_2$, $\ldots$, $x-\alpha_n$ annehmen, und setzen, da der Coëfficient der höchsten Potenz (nach \S~1) $=1$ ist,
\begin{equation}
\begin{aligned}
 f(x)&=(x-\alpha_1)(x-\alpha_2)\cdots(x-\alpha_n)\\
 &=x^n+a_1x^{n-1}+a_2x^{n-2}+\cdots+a_n.
\end{aligned}
\tag{1}
\end{equation}
Eine leichte Ueberlegung lässt folgendes Bildungsgesetz der Coëfficienten $a_1,a_2\ldots a_n$ erkennen:

Es ist $-a_1$ gleich der Summe der $\alpha$, $a_2$ gleich der Summe der Producte von je zweien der $\alpha$, $-a_3$ die Summe der Producte von je dreien der $\alpha$, allgemein $(-1)^\nu a_\nu$ die Summe der Producte von je $\nu$ der Grössen $\alpha$, oder in Formeln ausgedrückt:
\begin{equation}
\begin{aligned}
-a_1&=\sum \alpha_1,\\
+a_2&=\sum \alpha_1\alpha_2,\\
&\cdots\\
(-1)^\nu a_\nu&=\sum \alpha_1\alpha_2\cdots\alpha_\nu,\\
&\cdots\\
(-1)^n a_n&=\alpha_1\alpha_2\cdots\alpha_n.
\end{aligned}
\tag{2}
\end{equation}
Um aber noch deutlicher die Richtigkeit dieses Bildungsgesetzes einzusehen, bedient man sich der \gesperrt{vollständigen Induction}.

Man bestätigt die Richtigkeit zunächst in den ersten Fällen durch wirkliche Ausführung der Multiplication
%CONT>>
